\kafkasection{Calculate the mass-radius relation for white dwarfs, as follows: Use the exact expression for the Fermi momentum and the non-relativistic expression for pressure of a degenerate electron gas. Assume the white dwarf is pure ionized 12C. Integrate the equation of hydrostatic equilibrium for a constant density white dwarf to get the central pressure and hence the mass radius relation. Give your final result as R as a function of ($M/M_\solar$) with a numerical coefficient.}
The fermi momentum is:
\begin{equation}
    p_f = \ab(\frac{3h^3 n_e}{8\pi})^{1/3}
\end{equation}
This density of the electron is related to the nuclei density since our star is fully carbon-12
\begin{equation}
    n_e = \frac{Z}{A} \frac{\rho}{m_p} 
\end{equation}
Our fermi momentum is then:
\begin{equation}
    p_f = \ab(\frac{3h^3}{8\pi}\frac{Z\rho}{Am_p})^{1/3}
\end{equation}
We will use this in the pressure expression from momentum.
\begin{equation}
    P = \frac{8\pi}{15h^3 m_e}p_F^5 = \frac{8\pi}{15h^3m_e}\ab(\frac{3h^3}{8\pi}\frac{Z\rho}{Am_p})^{5/3}
\end{equation}
Where the hydrostatic equilibrium relation is as follows:
\begin{equation}
    \odv{P}{r} = -\rho \frac{Gm(r)}{r^2}
\end{equation}
We can then integrate this over $P$ for the LHS and over $r$ for the RHS. Our bounds of integration will be from the core (Where we have a central pressure $P_c$ and to the stellar radius where pressure is zero)
\begin{equation}
    P\evalat_{P_c}^0 = -G\rho \int_{0}^R \frac{m(r)}{r^2}\d r 
\end{equation}
Let's deal with this integral on the RHS:
\begin{equation}
    \int_0^R \frac{m(r)}{r^2}\d r = \int_0^R \frac{\rho (4/3)\pi r^3}{r^2}\d r = (2/3)\rho \pi R^2
\end{equation}
So the total RHS is then:
\begin{equation}
    -(2/3)G\rho^2 \pi R^2
\end{equation}
We then have:
\begin{equation}
    P_c = \frac{2G\rho^2\pi R^2}{3}
\end{equation}
Let's equate this to our pressure derived from fermi momentum:
\begin{equation}
    \frac{8\pi}{15h^3m_e}\ab(\frac{3h^3}{8\pi}\frac{Z\rho}{Am_p})^{5/3} = \frac{2G\rho^2\pi R^2}{3}
\end{equation}
Let's simplify math by considering that we have constant density (a requirement since we integrated over $r$) and that the density is the total mass $M$ divided by the volume.
\begin{equation}
    \frac{8\pi}{15h^3m_e} \ab(\frac{3h^3}{8\pi}\frac{Z}{Am_p} \frac{M}{(4/3)\pi R^3})^{5/3} = \frac{2G\pi R^2}{3}\frac{M^2}{(16/9)\pi^2 R^6}
\end{equation}
Now we begin the arduous journey of solving for $R$.
\begin{equation}
    \frac{8\pi h^5}{15h^3m_e R^5} \ab(\frac{9Z}{32\pi^2Am_p})^{5/3} M^{5/3} = \frac{3GM^2}{8\pi R^4}
\end{equation}
\begin{equation}
    R = \frac{64\pi^2h^2}{45Gm_e}\ab(\frac{9Z}{32\pi^2Am_p})^{5/3} M^{-1/3}
\end{equation}
So now we have $R$ as a function of $M$ with a constant out front. This is going to be hell to compute but it is (for carbon-12):
\begin{equation}
    R = \qty{3.60e19}{cm g^{1/3}} M^{-1/3}
\end{equation}
However, to get this in units of mass fraction of solar mass, we must divide by $M_\solar^{-1/3}$. This gives the constant as:
\begin{equation}
    \qty{3.60e19}{cm g^{1/3}} \cdot \frac{(\qty{1.989e33}{g})^{-1/3}}{M_\solar^{-1/3}} = \qty{2.86e8}{cm M_\solar^{-1/3}}
\end{equation}
So finally
\begin{equation}
    \boxed{R = \qty{2.86e8}{cm M_\solar^{1/3}} M^{-1/3}}
\end{equation}
Where $M$ is expressed as a fraction of solar mass.